\subsection{Визуальная навигация}

Зная координаты углов маркера на изображении, мы можем вычислить положение дрона относительно маркера. Это называется \textbf{Pose Estimation} (Оценка позы).

\subsubsection*{Определение положения (Pose Estimation)}
Для точного определения расстояния и угла наклона используется функция \texttt{cv2.aruco.estimatePoseSingleMarkers}.
Она возвращает два вектора:
\begin{itemize}
    \item \textbf{tvec (Translation Vector):} Вектор смещения $(x, y, z)$ — где маркер относительно камеры.
    \item \textbf{rvec (Rotation Vector):} Вектор вращения — как маркер наклонен.
\end{itemize}

Для работы этой функции необходимо знать параметры камеры (матрицу калибровки), которые можно получить с помощью скрипта калибровки.

\subsubsection*{Простая логика следования}
Даже без сложной математики можно реализовать следование за маркером:
\begin{enumerate}
    \item Найти центр маркера на изображении (среднее арифметическое углов).
    \item Найти центр кадра (обычно 320x240 для разрешения 640x480).
    \item Если $X_{marker} > X_{center}$, значит маркер справа $\rightarrow$ поворачиваем вправо.
    \item Если $X_{marker} < X_{center}$, значит маркер слева $\rightarrow$ поворачиваем влево.
\end{enumerate}

\begin{codefile}[python]{python/examples/12_follow_aruco.py}
\end{codefile}

Этот пример показывает базовую логику: дрон взлетает, висит в воздухе и ищет маркеры. Если он видит маркер с ID 42 (ответ на главный вопрос жизни, вселенной и всего такого), он совершает автоматическую посадку.

\begin{taskbox}[Посадка на базу]
Реализуйте алгоритм: дрон висит в воздухе. Если он видит маркер ID 1, он летит вперед. Если видит маркер ID 2, летит назад. Если видит маркер ID 0 — садится.
\end{taskbox}
